\chapter*{Introduction}
\addcontentsline{toc}{chapter}{Introduction}
\markboth{Introduction}{Introduction}

\dropcap{T}{he} problem of sequential investment under uncertainty motivates this thesis. Modern reinforcement learning offers tools for such decisions, but most methods collapse the return to its expectation. Distributional approaches instead model the full return distribution, as introduced by Dabney et~al.~\cite{dabney2018implicit} with Implicit Quantile Networks. Separately, epistemic uncertainty over a model's own estimates can be quantified through Evidential Deep Learning \cite{sensoy2018evidential}. This thesis combines both ideas \cite{dabney2018implicit, sensoy2018evidential} into a single risk-sensitive decision support system.

\section*{Box gallery (visual test)}

\begin{definition}[Invariants of unimodular symmetric forms]
 Let $Z$ be a finitely-generated free abelian group, and let $q : Z \times Z \to \mathbb{Z}$ be a unimodular form. To it we assign rank, type (parity), and signature.
 \end{definition}

% \begin{notation}
% In [MH73], unimodular symmetric forms on $Z$ are called inner products on $Z$. \\ 
% Forms of type I are odd and forms of type II are even.
% \end{notation}

% \begin{example}[$E_8$]
% Define the form $E_8$ on $\mathbb{Z}^8$. This form is unimodular, positive-definite and even, so $\operatorname{rank}(E_8) = \operatorname{sign}(E_8) = 8$.
% \end{example}

% \begin{theorem}[Freedman]
% Let $q$ be a unimodular, symmetric form. If $q$ is even, there is exactly one closed, oriented, simply-connected topological four-manifold whose intersection form is $q$.
% \end{theorem}

% \begin{proposition}[$w_2(M)$ is mod-2 reduction of characteristic vector]
% Let $M$ be a closed, connected, oriented four-manifold with intersection form $q_M$. \\
% Then for any characteristic vector $w \in H^2(M,\mathbb{Z})$, $w_2(M) = \mathrm{PD}(w) \bmod 2$.
% \end{proposition}

% \begin{corollary}[Closed four-manifolds admit Spin$^c$-structures]
% Every closed, connected, oriented four-manifold admits Spin$^c$-structures.
% \end{corollary}

% \begin{lemma}[Elkies]
% Let $q : Z \times Z \to \mathbb{Z}$ be a symmetric, unimodular bilinear form. If $q$ is not $\oplus(-1)$ nor $\oplus(+1)$, then there exists a characteristic vector $w$ with $|q(w,w)| < \operatorname{rank}(q)$.
% \end{lemma}

% \section*{Acknowledgements}
% First of all, I would like to thank the mathematical community whose documents make LaTeX typography worth studying. This section is included to test the visual spacing of front matter pages.